% Standalone square panel; set \MIPAView to 1 or 2 before input.
% Values: mapping_tau0025_flip/assignments.json, tau=0.025, normalized over10 prototypes.
% Manual illustration edits in the left panel: 6.1 -> 14.2, 33.5 -> 29.5, 11.5 -> 9.4, 8.6 -> 6.6.
\usetikzlibrary{calc}
\begingroup
\def\MIPAPrototypeData{%
  prototype-01/7.8/5.7,
  prototype-02/4.5/5.8,
  prototype-03/14.2/19.9,
  prototype-04/2.7/5.1,
  prototype-05/10.5/10.7,
  prototype-06/29.5/22.5,
  prototype-07/10.2/10.2,
  prototype-08/9.4/7.8,
  prototype-09/4.6/3.3,
  prototype-10/6.6/8.9%
}
\definecolor{MIPAReal}{RGB}{91,132,180}
\definecolor{MIPASynthetic}{RGB}{199,106,100}
\definecolor{MIPABlueFill}{RGB}{217,230,244}
\definecolor{MIPARedFill}{RGB}{248,222,218}
\definecolor{MIPAText}{RGB}{40,44,49}
\definecolor{MIPAMuted}{RGB}{91,101,111}
\colorlet{MIPAGrid}{MIPAMuted!17!white}
\colorlet{MIPAFrame}{MIPAMuted!28!white}
\ifnum\MIPAView=1
  \colorlet{MIPAAccent}{MIPAReal}
  \colorlet{MIPABarFill}{MIPABlueFill}
\else
  \colorlet{MIPAAccent}{MIPASynthetic}
  \colorlet{MIPABarFill}{MIPARedFill}
\fi
\edef\MIPAImageDirectory{images/view\MIPAView}
\begin{tikzpicture}[x=1bp,y=1bp,inner sep=0pt,outer sep=0pt,
  every node/.style={font=\rmfamily\upshape}]
  \path[use as bounding box] (-126,-126) rectangle (126,126);
  \fill[white] (-126,-126) rectangle (126,126);
  \foreach \MIPARadius in {60,86} {
    \draw[MIPAGrid,line width=.3bp,dash pattern=on .7bp off 2bp]
      (0,0) circle[radius=\MIPARadius bp];
  }
  \foreach \MIPAFile/\MIPAFirst/\MIPASecond
    [count=\MIPAIndex from 0] in \MIPAPrototypeData {
    \ifnum\MIPAView=1\let\MIPAProbability\MIPAFirst
    \else\let\MIPAProbability\MIPASecond\fi
    \pgfmathsetmacro{\MIPAAngle}{90-36*\MIPAIndex}
    \pgfmathsetmacro{\MIPAStart}{\MIPAAngle-12}
    \pgfmathsetmacro{\MIPAEnd}{\MIPAAngle+12}
    % Full colored bar lengths encode probabilities; numbers follow the outer caps.
    \pgfmathsetmacro{\MIPAOuter}{34+52*\MIPAProbability/35}
    \draw[MIPAGrid,line width=.3bp] (\MIPAAngle:24) -- (\MIPAAngle:88);
    \path[fill=MIPABarFill,draw=MIPAAccent!80!white,line width=.4bp,line join=round]
      (\MIPAStart:34) -- (\MIPAStart:\MIPAOuter)
      arc[start angle=\MIPAStart,end angle=\MIPAEnd,radius=\MIPAOuter bp]
      -- (\MIPAEnd:34)
      arc[start angle=\MIPAEnd,end angle=\MIPAStart,radius=34bp] -- cycle;
    \pgfmathsetmacro{\MIPALabelRadius}{\MIPAOuter+.6}
    % Keep the bottom of every number against its cap, including the lower half.
    \pgfmathsetmacro{\MIPALabelAngle}{\MIPAAngle-90}
    \node[text=MIPAText,inner sep=.5bp,rotate=\MIPALabelAngle,anchor=south,
      font=\rmfamily\bfseries\upshape\fontsize{9.5bp}{10bp}\selectfont]
      (MIPALabel-\MIPAIndex) at (\MIPAAngle:\MIPALabelRadius)
      {\MIPAProbability};
    \path let \p1=(MIPALabel-\MIPAIndex.east), \p2=(MIPALabel-\MIPAIndex.west),
      \p3=(MIPALabel-\MIPAIndex.north), \p4=(MIPALabel-\MIPAIndex.south),
      \n1={veclen(\x1-\x2,\y1-\y2)}, \n2={veclen(\x3-\x4,\y3-\y4)}
      in \pgfextra{\typeout{MIPA_LABEL:\MIPAView:\MIPAIndex:\n1:\n2:0pt:\MIPAOuter}
        \typeout{MIPA_BASE:\MIPAView:\MIPAIndex:\x4:\y4:\x3:\y3}};
  }
  \filldraw[fill=white,draw=MIPAMuted!24!white,line width=.4bp]
    (0,0) circle[radius=24bp];
  \draw[MIPAReal,line width=.55bp] (-14.7,-14.7) rectangle (14.7,14.7);
  \node at (0,0) {\includegraphics[width=28bp,height=28bp]{\MIPAImageDirectory/real.png}};
  \foreach \MIPAFile/\MIPAFirst/\MIPASecond
    [count=\MIPAIndex from 0] in \MIPAPrototypeData {
    \ifnum\MIPAView=1\let\MIPAProbability\MIPAFirst
    \else\let\MIPAProbability\MIPASecond\fi
    \pgfmathsetmacro{\MIPAAngle}{90-36*\MIPAIndex}
    \colorlet{MIPABorder}{MIPAFrame}\def\MIPABorderWidth{.4bp}
    \ifdim\MIPAProbability pt>20pt
      \colorlet{MIPABorder}{MIPAAccent}\def\MIPABorderWidth{.65bp}
    \fi
    \begin{scope}[shift={(\MIPAAngle:110)}]
      \draw[MIPABorder,line width=\MIPABorderWidth] (-14.7,-14.7) rectangle (14.7,14.7);
      \node at (0,0) {\includegraphics[width=28bp,height=28bp]{\MIPAImageDirectory/\MIPAFile.png}};
    \end{scope}
  }
\end{tikzpicture}
\endgroup
